\section{Legal Discussion}
\label{sec:legal}

\subsection{Federal Prohibition (\S104(e))}

The hypothetical Districting Integrity Act \S104(e) prohibits partisan data
as an input to federal congressional redistricting.
B.12 is not a proposal to violate this prohibition.
Rather, it is a mathematical analysis of what proportionality would cost if
the prohibition were relaxed --- and a characterisation of when it costs nothing
(free proportionality states) versus when it costs significantly.

The \S104(e) prohibition is a policy choice, not a mathematical necessity.
The choice has consequences: in expensive-proportionality states, geography-only
algorithms produce maps where 50\,\% of voters elect 35\,\% of representatives.
B.12 measures those consequences.

\subsection{State Legislative Applications}

State legislatures are not constrained by \S104(e).
Several state constitutions require proportional representation or prohibit
maps that unduly favour a political party \citep{lwv2018pa, harper2022nc}.

For state legislative redistricting, B.12 provides a rigorous operationalisation
of proportionality: set $\mathtt{ubvec}[1]$ to the minimum value that achieves
proportional seat allocation, and verify the Lorenz feasibility condition is
satisfied.
If Lorenz feasibility fails (expensive-proportionality state), the court can
be informed that geographic constraints make proportionality unachievable
with single-member districts --- a fact about geography, not about intent.

\subsection{The Evidentiary Role}

B.12's most immediate legal contribution is evidentiary.
For states where $\sigma$ is low (free or cheap proportionality), a court
can assess whether a geography-only algorithm's non-proportional outcome is
``natural'' (unavoidable given geography) or artificially produced.
If $\sigma < 0.05$ and the algorithm still produces a $-10$\,pp gap,
something other than geography is causing the shortfall.

For states where Lorenz feasibility is violated (expensive-proportionality),
the court has a rigorous basis for finding that proportionality is geometrically
impossible with contiguous single-member districts, which may support a
different remedy (e.g., multi-member districts).

\subsection{\emph{Rucho v.\ Common Cause} (2019)}

\emph{Rucho} held that federal courts cannot remedy partisan gerrymandering
under the federal Constitution.
B.12's analysis does not challenge this holding.
What B.12 adds is a precise characterisation of what ``partisan gerrymandering''
is on a spectrum: from zero partisan signal (geography-only, never gerrymandering)
to full partisan signal (explicitly partisan, clearly gerrymandering).
The \emph{Rucho} holding leaves open whether states can regulate intermediate
points on this spectrum --- and B.12 provides the mathematical tools to
characterise any point on the spectrum precisely.
